% =====================================================================
%  Mémoire de stage M2 IMD — Toky Nirina RAMAHERISOA — 2025/2026
%  VERSION DE TRAVAIL v1
%   - Sections 1 à 4 : rédaction réelle (lot 1)
%   - Sections 5 à 7 : encore au stade plan (\descr), rédigées au lot 2
%
%  Compilation : pdflatex (Overleaf), 2 passes.
%  Les balises \todo{...} (rouge) marquent les chiffres / captures à fournir.
% =====================================================================
\documentclass[11pt,a4paper]{article}

% ---------- Encodage / langue ----------
\usepackage[utf8]{inputenc}
\usepackage[T1]{fontenc}
\usepackage{lmodern}
\usepackage[french]{babel}

% ---------- Mise en page ----------
\usepackage[margin=2.4cm]{geometry}
\usepackage{setspace}
\usepackage{parskip}

% ---------- Maths / tableaux / figures ----------
\usepackage{amsmath,amssymb}
\usepackage{amsthm}
\usepackage{booktabs}
\usepackage{graphicx}
\usepackage{enumitem}
\usepackage{xcolor}
\usepackage{tikz}

% ---------- Liens ----------
\usepackage[hidelinks]{hyperref}
\hypersetup{
  pdftitle={Memoire M2 IMD - Generation PPC de non-benzenoides 5/6/7},
  pdfauthor={Toky Nirina Ramaherisoa}
}

% =====================================================================
%  Macros
% =====================================================================
\definecolor{descrgray}{gray}{0.35}
\newcommand{\descr}[1]{%
  {\small\itshape\color{descrgray}\textbf{Contenu prévu :} #1\par}\vspace{0.4em}}
\newcommand{\todo}[1]{{\bfseries\color{red}[À COMPLÉTER : #1]}}
\newcommand{\tbd}{{\bfseries\color{red}--?--}}  % cellule de tableau à remplir

\newcommand{\GD}{\ensuremath{G_D}}                 % graphe dual
\newcommand{\Tn}{\ensuremath{\mathcal{T}}}         % table de voisinage
\newcommand{\xtb}{\textsc{xTB}}
\newcommand{\ace}{\textsc{ace}}
\newcommand{\Aut}{\operatorname{Aut}}

\newtheorem{definition}{Définition}
\newtheorem{remarque}{Remarque}

\begin{document}

% =====================================================================
%  PAGE DE GARDE
% =====================================================================
\begin{titlepage}
\centering
% À FAIRE sur Overleaf : déposer amu.png, lis.png, ism2.png puis décommenter
% \includegraphics[height=1.7cm]{amu.png}\hfill
% \includegraphics[height=1.7cm]{lis.png}\hfill
% \includegraphics[height=1.7cm]{ism2.png}
{\large\bfseries Aix-Marseille Université\quad$\cdot$\quad LIS\quad$\cdot$\quad iSm2}\\[0.3em]
{\footnotesize\itshape [emplacement des logos — à ajouter sur Overleaf]}

\vspace{1.2cm}
{\large Mémoire de stage de Master 2}\\[0.2em]
{\large Informatique et Mathématiques Discrètes (IMD)}\\[0.2em]
{\large Année universitaire 2025--2026}

\vspace{1.6cm}
\rule{\linewidth}{0.5pt}\\[0.6em]
{\LARGE\bfseries Génération automatique et exhaustive de structures\\[0.25em]
moléculaires bi-dimensionnelles par programmation\\[0.25em]
par contraintes\par}
\vspace{0.5em}
{\large focus sur les non-benzénoïdes à cycles 5, 6 et 7 atomes\\
et validation quantique semi-empirique\par}
\rule{\linewidth}{0.5pt}

\vspace{1.6cm}
{\large\bfseries Toky Nirina RAMAHERISOA}

\vspace{1.4cm}
\begin{tabular}{rl}
\textbf{Encadrants} & Nicolas \textsc{Prcovic} et Cyril \textsc{Terrioux}\\
                    & \emph{LIS, équipe COALA, Aix-Marseille Université}\\[0.6em]
\textbf{Partenaires} & Yannick \textsc{Carissan} et Denis \textsc{Hagebaum-Reignier}\\
                     & \emph{Institut des Sciences Moléculaires de Marseille (iSm2)}\\
\end{tabular}

\vfill
{\small Document de travail --- version v1}
\end{titlepage}

\tableofcontents
\newpage

% =====================================================================
\section{Introduction}
% =====================================================================

\subsection{Contexte}

Les \emph{hydrocarbures aromatiques polycycliques} (HAP) sont des molécules
constituées uniquement de carbone et d'hydrogène, dont les atomes de carbone
forment des cycles fusionnés. Ils sont étudiés dans de nombreux domaines ---
science des matériaux, nano-électronique moléculaire, opto-électronique,
astrochimie --- pour leurs propriétés électroniques, magnétiques et
spectroscopiques remarquables. Parmi les HAP, les \emph{benzénoïdes} forment la
sous-famille composée exclusivement d'hexagones fusionnés par une arête ;
le benzène, le naphtalène ou le coronène en sont des représentants classiques.
Cette famille a fait l'objet de nombreux travaux, et notamment de la thèse de
A.~Varet~\cite{varet2022}, qui a montré que la \emph{programmation par
contraintes} (PPC) est un outil particulièrement adapté à la génération
exhaustive de benzénoïdes répondant à des propriétés données.

L'étude des HAP comportant des cycles à \emph{cinq} ou \emph{sept} atomes ---
on parle alors de structures \emph{non-benzénoïdes} --- est plus récente.
L'azulène (un pentagone et un heptagone accolés) en est l'exemple emblématique.
L'introduction de tels cycles modifie en profondeur la chimie de la molécule :
elle brise la régularité du pavage hexagonal, induit une courbure locale, et
favorise l'apparition de sites radicalaires. Ces structures intéressent donc
les chimistes, mais leur génération systématique reste largement inexplorée.

\subsection{Objectif et problématique}

L'objectif du stage est le suivant : \emph{étant donné un benzénoïde, énumérer
de manière exhaustive toutes les substitutions admissibles dans lesquelles
chaque hexagone est soit conservé, soit remplacé par un cycle à $5$ ou $7$
atomes}. La difficulté est double, à la fois combinatoire et géométrique. Nous
en donnons l'énoncé précis et l'analyse détaillée à la section~\ref{sec:probleme},
une fois introduit le vocabulaire nécessaire.

\subsection{Positionnement et contributions}

Ce stage se place dans la continuité directe de la thèse de Varet~\cite{varet2022} :
nous reprenons son cadre (modélisation PPC sur le \emph{graphe dual} d'un
benzénoïde, notions d'aromaticité) et nous l'étendons au cas, non traité
jusqu'ici, des cycles à $5$ et $7$ atomes. Nos contributions sont les suivantes.

\begin{itemize}[leftmargin=1.4em]
  \item Un \textbf{modèle PPC} qui énumère exhaustivement les substitutions
    $5/6/7$ d'un benzénoïde, fondé sur une \emph{table de voisinage} encodant
    les configurations locales chimiquement réalisables.
  \item Un \textbf{pipeline de validation physique} qui reconstruit en trois
    dimensions chaque solution du modèle, l'optimise par une méthode quantique
    semi-empirique (\xtb{}), puis juge sa planéité.
  \item Une \textbf{étude expérimentale systématique} de contraintes
    additionnelles visant à augmenter la proportion de molécules planes
    produites. Ce point mérite une précision : la planéité d'une molécule est
    une propriété \emph{globale} qui ne se laisse pas exprimer entièrement par
    des contraintes locales --- c'est une difficulté intrinsèque au problème, et
    non un défaut de modélisation. Notre démarche consiste à identifier
    expérimentalement les contraintes qui, sans capturer la planéité au sens
    strict, en améliorent significativement le rendement.
  \item Une \textbf{extension de l'analyse électronique} (structures de Kekulé,
    couvertures de Clar, \emph{Ring Bond Order}) aux cycles à $5$ et $7$ atomes,
    là où les outils existants se limitent aux hexagones.
\end{itemize}

Le mémoire est organisé comme suit. La section~\ref{sec:prelim} introduit les
deux corpus de notions nécessaires --- la PPC et la chimie des HAP. La
section~\ref{sec:probleme} détaille le problème et ses difficultés. La
section~\ref{sec:modele} présente notre modélisation. La
section~\ref{sec:realisation} décrit la réalisation (reconstruction, validation,
analyse). La section~\ref{sec:xp} rapporte l'étude expérimentale, et la
section~\ref{sec:conclusion} conclut.

% =====================================================================
\section{Préliminaires et état de l'art}
\label{sec:prelim}
% =====================================================================

Notre travail se situe à l'interface de l'informatique et de la chimie. Cette
section donne au lecteur les notions nécessaires des deux côtés : d'abord le
paradigme de la programmation par contraintes, ensuite le vocabulaire de chimie
employé dans tout le mémoire.

\subsection{Programmation par contraintes (PPC)}

La \emph{programmation par contraintes} est un paradigme déclaratif de
résolution de problèmes combinatoires. Plutôt que de décrire \emph{comment}
résoudre un problème, on se contente d'en \emph{décrire les solutions}, puis on
confie la recherche à un solveur générique.

\begin{definition}[Réseau de contraintes / CSP]
Un \emph{problème de satisfaction de contraintes} est un triplet
$(X, D, C)$ où $X = \{x_1, \dots, x_n\}$ est un ensemble de \emph{variables},
$D = \{D_1, \dots, D_n\}$ associe à chaque variable $x_i$ un \emph{domaine}
fini $D_i$ de valeurs possibles, et $C$ est un ensemble de \emph{contraintes}.
Chaque contrainte porte sur un sous-ensemble de variables et restreint les
combinaisons de valeurs qu'elles peuvent prendre simultanément. Une
\emph{solution} est une affectation de chaque variable à une valeur de son
domaine satisfaisant toutes les contraintes.
\end{definition}

La résolution combine une \emph{recherche arborescente} (on affecte les
variables une à une, en revenant en arrière en cas d'échec) et de la
\emph{propagation de contraintes} (à chaque étape, on retire des domaines les
valeurs qui ne peuvent appartenir à aucune solution, par exemple par
\emph{cohérence d'arc}). L'intérêt majeur du paradigme est qu'il n'est pas
nécessaire d'écrire un algorithme spécifique à chaque problème : il suffit de
formuler le modèle, puis d'invoquer un solveur traité comme une \emph{boîte
noire}. Dans ce travail, nous écrivons les modèles avec la bibliothèque Python
\textsc{PyCSP3} et nous les résolvons avec le solveur \ace{} (\emph{Abscon
Constraint Engine})~\cite{pycsp3ace}, qui sait énumérer l'ensemble des solutions
d'une instance.

Deux familles de contraintes nous seront utiles. Une \textbf{contrainte
extensionnelle} (ou \emph{de table}) liste explicitement les combinaisons de
valeurs autorisées pour un groupe de variables ; c'est la forme naturelle pour
encoder un catalogue de configurations admissibles. Une \textbf{contrainte
globale} capture un motif récurrent sur un nombre arbitraire de variables : par
exemple $\mathit{Sum}$ (une somme bornée) ou $\mathit{LexIncreasing}$ (un ordre
lexicographique entre deux vecteurs de variables), que nous emploierons pour la
\emph{rupture de symétrie}. Enfin, la souplesse de la PPC --- activer, désactiver
ou ajouter une contrainte, puis ré-énumérer rapidement les solutions --- est un
atout décisif pour un travail mené en dialogue continu avec les chimistes.

\subsection{Notions de chimie}

Nous regroupons ici les termes de chimie employés dans le mémoire. Le lecteur
informaticien peut les lire comme des définitions sur des graphes ; le lecteur
chimiste y reconnaîtra le vocabulaire usuel.

\subsubsection{HAP, benzénoïdes et graphe dual}

Un \textbf{benzénoïde} est un assemblage d'hexagones réguliers du pavage
hexagonal, fusionnés par des arêtes. On peut le voir de deux façons. Le
\emph{graphe carbone} a pour sommets les atomes de carbone et pour arêtes les
liaisons carbone--carbone. Le \emph{graphe dual} \GD{}, central dans tout ce
travail, résume la combinatoire des hexagones.

\begin{definition}[Graphe dual]
Soit $B$ un benzénoïde à $h$ hexagones. Son graphe dual $\GD = (V_D, E_D)$ a
pour sommets $V_D = \{0, \dots, h-1\}$ (un sommet par hexagone), et
$(u,v) \in E_D$ si et seulement si les hexagones $u$ et $v$ partagent une arête
carbone--carbone.
\end{definition}

On appelle \textbf{bord} de \GD{} l'ensemble des hexagones de degré strictement
inférieur à $6$, c'est-à-dire ceux qui ne sont pas entièrement entourés de
voisins. La figure~\ref{fig:dual} illustre cette construction sur un benzénoïde
linéaire à trois hexagones.

\begin{figure}[h]
\centering
\begin{tikzpicture}[scale=0.62]
  \foreach \k/\lab in {0/v_0, 1/v_1, 2/v_2} {
    \begin{scope}[xshift={\k*2*1.2*cos(30)*1cm}]
      \foreach \i in {0,...,5} {
        \coordinate (P\i) at ({cos(60*\i+30)*1.2}, {sin(60*\i+30)*1.2});
      }
      \draw[thick] (P0)--(P1)--(P2)--(P3)--(P4)--(P5)--cycle;
      \node at (0,0) {$\lab$};
    \end{scope}
  }
\end{tikzpicture}
\hspace{1.2cm}
\begin{tikzpicture}[scale=1]
  \node[circle,draw,inner sep=2pt] (a) at (0,0) {$0$};
  \node[circle,draw,inner sep=2pt] (b) at (1.4,0) {$1$};
  \node[circle,draw,inner sep=2pt] (c) at (2.8,0) {$2$};
  \draw (a)--(b)--(c);
\end{tikzpicture}
\caption{Un benzénoïde linéaire à trois hexagones (gauche) et son graphe dual
\GD{} (droite) : un chemin à trois sommets.}
\label{fig:dual}
\end{figure}

\subsubsection{Aromaticité et délocalisation}

L'\textbf{aromaticité} traduit la stabilité particulière de certains cycles
conjugués. Dans un hexagone de type benzène, chaque carbone porte un électron
$\pi$ qui n'est pas localisé sur une liaison précise : les électrons sont
\emph{délocalisés} et circulent autour du cycle. Plus une molécule admet de
représentations équivalentes de ses liaisons doubles, plus la délocalisation est
étendue, et plus la molécule est stable.

\subsubsection{Structures de Kekulé}

Une \textbf{structure de Kekulé} est un étiquetage des liaisons carbone--carbone
en simples et doubles tel que chaque carbone porte exactement une double
liaison. En termes de graphes, c'est un \emph{couplage parfait} du graphe
carbone. Une molécule possède en général plusieurs structures de Kekulé.

\subsubsection{Sextets et couvertures de Clar}

Un \textbf{sextet de Clar} (ou rond de Clar) est un hexagone dont les six
électrons $\pi$ forment un cycle aromatique autonome. Une \textbf{couverture de
Clar} place des sextets sur un ensemble d'hexagones deux à deux disjoints, et le
\textbf{nombre de Clar} est le nombre maximal de sextets que l'on peut placer.
La règle de Hückel ($4n+2$ électrons $\pi$, soit $6$ pour $n=1$) impose qu'un
sextet occupe exactement six atomes : seuls les \emph{hexagones} peuvent porter
un sextet, jamais un pentagone ni un heptagone.

\subsubsection{Ring Bond Order (RBO)}

L'\textbf{ordre de liaison} d'une arête est la fraction des structures de Kekulé
dans lesquelles cette arête est double ; c'est une mesure de la délocalisation
locale. Le \textbf{Ring Bond Order} d'un cycle est la somme des ordres de
liaison de ses arêtes. Il quantifie à quel point un cycle « ressemble » à un
cycle aromatique sur l'ensemble des Kekulé.

\subsubsection{Sites radicalaires}

Lorsqu'une molécule n'admet aucun couplage parfait (par exemple à cause d'un
nombre impair d'atomes ou d'une topologie particulière), certains carbones
restent non couverts : ils portent un \textbf{électron $\pi$ non apparié}, ou
\emph{radical}. La parité impaire des cycles à $5$ et $7$ atomes favorise
précisément l'apparition de tels sites, ce qui rend les structures
non-benzénoïdes intéressantes pour la chimie radicalaire.

\subsubsection{Planéité et courbure}

Une molécule est \textbf{plane} si tous ses atomes se placent (à l'équilibre
énergétique) dans un même plan. Pour un pavage d'hexagones réguliers, la somme
des angles autour de chaque sommet vaut $360^\circ$ et la structure est plane.
Un pentagone introduit localement un \emph{défaut d'angle} positif (la matière
tend à se refermer, comme dans un fullerène) et un heptagone un défaut négatif
(la matière tend à onduler). Cette intuition se formalise par le théorème de
Gauss--Bonnet discret : la \emph{courbure} cumulée localement vaut
$\tfrac{\pi}{3}(n_5 - n_7)$, où $n_5$ et $n_7$ comptent pentagones et heptagones
au voisinage. Un pentagone et un heptagone adjacents (motif \emph{Stone--Wales})
compensent mutuellement leur courbure et préservent la planéité ; un défaut
isolé, au contraire, tord la molécule.

\subsubsection{Validation quantique semi-empirique (\xtb{})}

Pour juger si une molécule est réellement plane, on cherche sa géométrie
d'\emph{équilibre} : la disposition des atomes qui minimise l'énergie. Les
méthodes \emph{ab initio} (DFT) sont précises mais coûteuses. Nous utilisons une
méthode \emph{semi-empirique}, GFN2-\xtb{}~\cite{xtb}, qui approxime l'énergie à
un coût bien moindre et fournit une géométrie optimisée exploitable à grande
échelle. C'est le sens de la « validation quantique semi-empirique » du titre.

% =====================================================================
\section{Le problème}
\label{sec:probleme}
% =====================================================================

Maintenant que le vocabulaire est posé, nous pouvons énoncer précisément le
problème et analyser ses deux sources de difficulté.

\subsection{Énoncé : substitutions $5/6/7$ d'un benzénoïde}

On part d'un benzénoïde $B$ à $h$ hexagones, représenté par son graphe dual
\GD{}. Une \textbf{substitution} affecte à chaque hexagone $v$ une taille de
cycle $x_v \in \{5, 6, 7\}$ : conserver l'hexagone ($x_v = 6$), le contracter en
pentagone ($x_v = 5$) ou le dilater en heptagone ($x_v = 7$). La substitution
doit respecter la \emph{topologie} de $B$ (les arêtes partagées entre cycles
voisins sont préservées) et conserver le nombre total d'atomes de carbone. On
cherche à \textbf{énumérer toutes les substitutions admissibles}, c'est-à-dire
toutes celles qui correspondent à une molécule chimiquement et géométriquement
réalisable, sans doublon (deux substitutions identiques à une symétrie près ne
comptent que pour une).

La figure~\ref{fig:subst} illustre les trois opérations élémentaires sur un
cycle isolé.

\begin{figure}[h]
\centering
\begin{tikzpicture}[scale=0.7]
  % pentagone
  \begin{scope}[xshift=-4.2cm]
    \foreach \i in {0,...,4} {\coordinate (A\i) at ({cos(72*\i+90)*1.1},{sin(72*\i+90)*1.1});}
    \draw[thick] (A0)--(A1)--(A2)--(A3)--(A4)--cycle;
    \node at (0,-1.7) {pentagone (5)};
  \end{scope}
  % hexagone
  \foreach \i in {0,...,5} {\coordinate (B\i) at ({cos(60*\i+30)*1.2},{sin(60*\i+30)*1.2});}
  \draw[thick] (B0)--(B1)--(B2)--(B3)--(B4)--(B5)--cycle;
  \node at (0,-1.7) {hexagone (6)};
  % heptagone
  \begin{scope}[xshift=4.2cm]
    \foreach \i in {0,...,6} {\coordinate (C\i) at ({cos(360/7*\i+90)*1.3},{sin(360/7*\i+90)*1.3});}
    \draw[thick] (C0)--(C1)--(C2)--(C3)--(C4)--(C5)--(C6)--cycle;
    \node at (0,-1.7) {heptagone (7)};
  \end{scope}
  \draw[->,thick] (-2.7,0)--(-1.6,0);
  \draw[->,thick] (1.6,0)--(2.7,0);
\end{tikzpicture}
\caption{Les trois tailles de cycle d'une substitution. Contracter un hexagone
en pentagone retire un sommet, le dilater en heptagone en ajoute un.}
\label{fig:subst}
\end{figure}

\subsection{La difficulté combinatoire}

Avec trois valeurs possibles par hexagone, un benzénoïde à $h$ hexagones admet a
priori $3^h$ substitutions. Ce nombre croît exponentiellement, et la grande
majorité de ces affectations ne correspondent à \emph{aucune} molécule valide :
elles violent la conservation du carbone, ou décrivent des voisinages
localement impossibles à reconstruire. Énumérer naïvement puis filtrer est donc
hors de portée dès que $h$ dépasse quelques unités. C'est exactement le type de
problème fortement combinatoire pour lequel la PPC est conçue : on encode les
conditions d'admissibilité comme des contraintes, et le solveur n'explore que
les affectations cohérentes.

\subsection{La difficulté géométrique et topologique}

La seconde difficulté est de nature différente, et c'est la plus profonde. Dans
un benzénoïde, tous les cycles sont des hexagones réguliers d'angle interne
$120^\circ$, et le pavage est plan par construction. Dès que l'on introduit des
pentagones (angle interne $108^\circ$) et des heptagones (environ $128{,}6^\circ$),
plusieurs grandeurs cessent d'être fixées :

\begin{itemize}[leftmargin=1.4em]
  \item les \textbf{angles} aux sommets ne se complètent plus à $360^\circ$, ce
    qui crée une courbure locale ;
  \item les \textbf{distances} inter-atomiques et les longueurs de liaison
    s'ajustent pour relâcher les tensions ;
  \item la \textbf{topologie} de fusion impose que les cycles voisins partagent
    exactement une arête, ce qui n'est pas toujours réalisable selon la
    disposition des voisins ;
  \item la \textbf{planéité} qui en résulte n'est pas garantie : la molécule peut
    se courber ou onduler pour minimiser son énergie.
\end{itemize}

Le point crucial est que la planéité est une propriété \emph{globale} : elle
dépend de l'agencement de tous les défauts à la fois, et ne se déduit pas
seulement des voisinages locaux. On peut filtrer localement les configurations
manifestement irréalisables, mais on ne peut pas, par des contraintes locales
seules, garantir qu'une molécule entière sera plane. Cette limite est
\emph{intrinsèque au problème}. Elle motive l'architecture en deux temps de notre
travail : un modèle PPC qui énumère les candidats en appliquant les conditions
\emph{exprimables} (section~\ref{sec:modele}), suivi d'une validation physique
qui tranche la planéité \emph{réelle} de chaque candidat
(section~\ref{sec:realisation}). Et elle motive l'étude expérimentale
(section~\ref{sec:xp}), qui cherche les contraintes améliorant au mieux le
rendement de planéité.

% =====================================================================
\section{Modélisation}
\label{sec:modele}
% =====================================================================

Nous présentons à présent le modèle PPC qui énumère les substitutions
admissibles. Il comporte une variable par hexagone et un ensemble de contraintes
\emph{structurelles}, toujours actives, que nous décrivons ici. Les contraintes
\emph{additionnelles}, qui font l'objet de l'étude expérimentale, sont
introduites à la section~\ref{sec:xp}.

\subsection{Variables et domaines}

À chaque hexagone $v$ du graphe dual on associe une variable $x_v$ donnant la
taille du cycle substitué :
\[
  x_v \in \{5, 6, 7\}, \qquad v \in V_D.
\]
Avant la résolution, un \emph{prétraitement} réduit le domaine de certaines
variables en exploitant la structure de \GD{}. Pour chaque hexagone $v$, on
définit son \textbf{motif d'occupation} $\sigma(v) = (\sigma_0, \dots, \sigma_5)
\in \{0,1\}^6$, où $\sigma_i = 1$ si le $i$-ème côté de $v$ est partagé avec un
voisin, et $0$ s'il est libre. On note $b(v)$ le nombre de \emph{blocs} de côtés
libres consécutifs (sur le cycle des six côtés).

Un hexagone est dit \textbf{gelé} --- forcé à rester un hexagone, domaine réduit
à $\{6\}$ --- dans deux cas :
\begin{itemize}[leftmargin=1.4em]
  \item $\deg_{\GD}(v) = 6$ : l'hexagone est entouré de six voisins, n'a aucun
    côté libre, et ne peut donc ni se contracter ni se dilater ;
  \item $b(v) \ge 2$ : ses côtés libres sont fragmentés en plusieurs blocs ;
    transformer le cycle exigerait de déformer un bloc indépendamment des
    autres, ce qui est topologiquement impossible.
\end{itemize}
Les autres hexagones, dits \emph{libres}, conservent le domaine $\{5, 6, 7\}$.

\subsection{Contrainte C1 --- Conservation du carbone}

Le nombre total d'atomes de carbone doit être préservé par la substitution.
Comme un pentagone retire un carbone et un heptagone en ajoute un (relativement
à l'hexagone), cela s'écrit
\[
  \sum_{v \in V_D} x_v = 6h,
\]
condition équivalente à $n_5 = n_7$ : il y a autant de pentagones que
d'heptagones. La solution triviale « tous les hexagones » (le benzénoïde de
départ) satisfait C1 ; comme nous cherchons les substitutions
\emph{non-benzénoïdes}, elle est écartée de l'énumération.

\subsection{Contrainte C3 --- Voisinage admissible et table de voisinage}

C'est le cœur du filtrage chimique. Pour chaque hexagone libre $v$, on impose
que la configuration locale formée par $v$ et ses voisins soit \emph{réalisable}.
Cette condition est encodée par une \textbf{contrainte extensionnelle} : le
tuple ordonné
\[
  \tau_v = (x_v, x_{u_1}, x_{u_2}, \dots, x_{u_k}),
\]
où $u_1, \dots, u_k$ sont les voisins de $v$ dans \GD{} pris dans l'ordre
cyclique, doit appartenir à une \textbf{table de voisinage} \Tn{} pré-calculée.

Un point essentiel distingue notre table de celle des benzénoïdes purs : elle
est \emph{sensible à la disposition} des voisins. Deux configurations ayant les
mêmes tailles de voisins mais disposées différemment autour du cycle central ne
sont pas équivalentes, car elles ne se reconstruisent pas de la même manière en
trois dimensions. La position des voisins est donc encodée explicitement dans
les tuples de la table (un $0$ marquant un côté libre).

La table est \textbf{construite hors-ligne} : on énumère les configurations
locales candidates (un cycle central entouré de voisins de tailles données dans
une disposition donnée), on reconstruit chacune en 3D, on l'optimise par \xtb{},
et on ne retient que celles qui sont planes (au sens du critère décrit en
section~\ref{sec:realisation}). Ce catalogue est \emph{paramétrable} : il pourra
être régénéré ou affiné lorsque de nouvelles études chimiques préciseront les
critères d'admissibilité. La version utilisée dans ce travail comporte
\textbf{721} entrées, réparties selon la taille du cycle central :
\[
  |\mathcal{T}(5)| = 63, \qquad |\mathcal{T}(6)| = 291, \qquad |\mathcal{T}(7)| = 367.
\]

\begin{remarque}
Un hexagone de degré $1$ (un seul voisin) n'a pas besoin de contrainte de table :
un cycle accolé à un unique voisin par une arête est toujours réalisable, quelle
que soit sa taille.
\end{remarque}

\subsection{Rupture de symétrie}

Deux substitutions qui ne diffèrent que par une symétrie du benzénoïde décrivent
la \emph{même} molécule. Pour ne les compter qu'une fois, on exploite le groupe
des automorphismes du graphe dual, $\Aut(\GD)$ (les permutations des sommets qui
préservent les arêtes), dont on calcule un ensemble générateur
$\{\pi_1, \dots, \pi_m\}$. Pour chaque générateur $\pi$, on impose une contrainte
\emph{lex-leader}
\[
  (x_0, x_1, \dots, x_{h-1}) \;\le_{\text{lex}}\; (x_{\pi(0)}, x_{\pi(1)}, \dots, x_{\pi(h-1)}),
\]
qui ne retient, dans chaque orbite de solutions équivalentes, que le représentant
minimal pour l'ordre lexicographique. Cette technique élimine les doublons et
réduit fortement l'espace de recherche. Sa validité repose sur l'unicité du
plongement planaire des graphes considérés (théorème de Whitney~\cite{whitney}
et 2-connexité des benzénoïdes).

\subsection{Résolution}

Le modèle (variables, domaines, contraintes C1, C3 et lex-leader) est écrit en
\textsc{PyCSP3}, puis compilé au format standard XCSP3. Le solveur \ace{} est
ensuite invoqué pour \textbf{énumérer toutes les solutions} de l'instance.
Chaque solution est un vecteur de tailles $(x_0, \dots, x_{h-1})$ qui sera, à
l'étape suivante, reconstruit en une molécule tridimensionnelle puis validé.

% =====================================================================
\section{Réalisation : des solutions du modèle aux molécules validées et analysées}
\label{sec:realisation}

Le modèle de la section~\ref{sec:modele} produit des \emph{vecteurs de tailles}
$(x_0, \dots, x_{h-1})$ : des objets abstraits. Cette section décrit comment on
en fait des molécules réelles, comment on tranche leur planéité, et comment on
analyse leurs propriétés électroniques. Elle regroupe l'essentiel du travail
d'implémentation.

\subsection{Reconstruction 3D des solutions}

Chaque solution est transformée en une molécule tridimensionnelle (carbones et
hydrogènes, avec coordonnées) en trois phases.

\begin{enumerate}[leftmargin=1.4em]
  \item \textbf{Modifications topologiques.} Pour chaque hexagone dont la taille
    assignée diffère de $6$, on modifie le cycle : on retire un sommet pour
    obtenir un pentagone, on en ajoute un pour un heptagone. Ces modifications
    ne portent que sur les côtés \emph{libres} ($\sigma_i = 0$), afin de
    préserver les arêtes partagées avec les voisins.
  \item \textbf{Placement géométrique.} Un parcours en largeur (BFS) du graphe
    dual place les cycles comme des polygones réguliers fusionnés (longueur de
    liaison carbone--carbone $\approx 1{,}4$~\AA). On part d'un cycle ancré (le
    plus haut degré), puis chaque cycle enfant est positionné par rapport à son
    parent via l'arête partagée, du bon côté pour que la structure ne se replie
    pas sur elle-même.
  \item \textbf{Assemblage moléculaire.} Les coordonnées sont injectées dans un
    graphe moléculaire ; les doubles liaisons sont placées par un couplage de
    cardinalité maximale (algorithme de Blossom d'Edmonds~\cite{edmonds1965}) ;
    les hydrogènes sont ajoutés sur les carbones de bord (longueur
    $\approx 1{,}088$~\AA) ; le tout est exporté au format XYZ.
\end{enumerate}

La géométrie ainsi produite est \emph{plane par construction} (tous les atomes
dans le plan $z=0$). C'est l'étape de validation qui révélera si la molécule est
réellement plane à l'équilibre.

\subsection{Validation par \xtb{} : protocole de dynamique moléculaire}

L'optimisation \xtb{} procède par descente de gradient. Or si la géométrie
initiale est \emph{parfaitement} plane, les forces perpendiculaires au plan sont
nulles par symétrie : l'optimiseur reste alors piégé dans un \textbf{minimum
plat parasite} et conclut à tort que la molécule est plane. Nous l'avons
constaté en comparant deux géométries initiales planes du même benzénoïde, qui
convergeaient l'une vers une forme plane, l'autre vers une forme courbée à
$58{,}6^\circ$.

Le protocole retenu contourne ce piège par une \textbf{dynamique moléculaire}
courte (GFN2-\xtb{}, $\sim 300$~K) qui « secoue » la molécule pour la sortir du
faux équilibre, suivie d'une optimisation (\texttt{--opt tight}) depuis la
dernière configuration, qui fournit la géométrie d'équilibre. Un filtre
anti-fragmentation rejette les trajectoires où un atome s'éjecte. C'est
l'équivalent numérique de la procédure suggérée par le chimiste : \emph{déplacer
un atome et vérifier s'il revient à sa place}. Deux variantes ont été testées :
\texttt{multi-runs} (plusieurs optimisations indépendantes avec une légère
perturbation aléatoire en $z$, la molécule étant déclarée plane si \emph{au
moins une} converge plane) et \texttt{z-perturb} (perturbation déterministe,
calculée par hachage du fichier, donc reproductible). Le protocole MD est la
stratégie par défaut.

\subsection{Test de planéité par analyse en composantes principales}

Sur la géométrie optimisée, on teste la planéité par une \textbf{analyse en
composantes principales} (ACP) des positions des carbones. On centre les
coordonnées, on diagonalise leur matrice de covariance $3\times3$, et l'on
obtient le \emph{plan moyen} (engendré par les deux premières composantes) et sa
normale $v_3$. On en déduit plusieurs métriques : l'angle entre $v_3$ et l'axe
$Z$ global, le \emph{rmsd} des distances des atomes au plan, la hauteur
(épaisseur) de la molécule, et surtout \texttt{max\_angle\_deg}, qui mesure les
bombements \emph{locaux} cycle par cycle. Le critère de décision, validé par le
chimiste, est
\[
  \text{molécule plane} \iff \texttt{max\_angle\_deg} \le 10^\circ .
\]
En combinant le statut de la MD et ce test, chaque solution reçoit un verdict :
\textsc{plan}, \textsc{non-plan}, ou \textsc{échec MD}.

\subsection{Analyse électronique des structures}

Une fois les structures planes obtenues, nous avons développé un outil pour en
étudier les propriétés électroniques. Ces analyses sont \emph{statiques} : elles
opèrent directement sur le graphe carbone (sans nouveau calcul quantique) et
sont donc quasi instantanées.

\begin{itemize}[leftmargin=1.4em]
  \item \textbf{Kekulé et radicaux.} On calcule un couplage de cardinalité
    maximale (Blossom) ; les Kekulé sont énumérés par retour arrière (avec un
    plafond). Les carbones non couverts sont les sites radicalaires.
  \item \textbf{Couvertures de Clar.} On énumère les sous-ensembles d'hexagones
    deux à deux disjoints (au plus $2^{n_{\text{hex}}}$, trivial pour
    $n_{\text{hex}} \le 9$) et l'on couple le résidu ; seuls les hexagones
    portent un sextet (règle de Hückel).
  \item \textbf{Extension du RBO aux cycles $5$ et $7$.} La définition de
    Pauling--Randić~\cite{randic}, reprise par Varet, ne concerne que les
    hexagones. Notre contribution est de conserver la formule
    $\mathrm{CBO}(C) = \sum_{e \in C} \mathrm{bo}(e)$ \emph{sommée sur tout
    cycle quelle que soit sa taille}, en affichant le ratio
    $\mathrm{CBO}/\mathrm{cbo_{max}}$ pour rester fidèle au contexte. Le cas
    radicalaire (aucune Kekulé stricte) n'est pas défini au sens de
    Pauling--Randić.
\end{itemize}

\begin{remarque}[Pourquoi l'énumération directe plutôt qu'une méthode polynomiale]
Le comptage rapide des Kekulé des benzénoïdes (méthode de
Rispoli~\cite{rispoli2001}) repose sur le caractère \emph{biparti} du graphe
carbone. Les cycles impairs ($5$ ou $7$) brisent cette bipartition ; la
généralisation passe par le théorème de Kasteleyn~\cite{kasteleyn1967}
(orientation et Pfaffien), d'un coût d'implémentation jugé hors du périmètre du
stage. Nous nous en tenons à l'énumération directe, suffisante car les cycles
$5/7$ contraignent fortement les couplages.
\end{remarque}

À ce stade, cet outillage est \emph{opérationnel} mais l'exploration reste
préliminaire : nous avons par exemple cherché un lien entre planéité et nombre
de radicaux sans dégager de corrélation exploitable. Nous le présentons comme un
moyen d'analyse des structures produites et une piste à approfondir.

\subsection{Mise en œuvre logicielle}

Le stage relevant de l'IMD et non du génie logiciel, nous décrivons ce volet
brièvement. L'ensemble du pipeline est intégré dans une application web (le
\emph{designer}) : un canevas hexagonal permet de dessiner ou d'importer un
benzénoïde et de choisir la configuration de contraintes, puis de lancer la
génération. Les résultats (solutions, géométries XYZ compressées, agrégats) sont
persistés dans une base SQLite, et les campagnes sur le grand corpus sont
distribuées sur un cluster. Un visualiseur 3D propose quatre modes d'affichage
(défaut, Kekulé, Clar, RBO). Ce travail d'industrialisation garantit surtout la
\emph{reproductibilité} des expériences décrites ci-dessous.

% =====================================================================
\section{Étude expérimentale : quelles contraintes améliorent la planéité ?}
\label{sec:xp}
% =====================================================================

Comme établi en section~\ref{sec:probleme}, la planéité est une propriété
globale qu'aucun jeu de contraintes locales ne capture exactement. La question
pratique devient alors : \emph{quelles contraintes additionnelles augmentent le
plus la proportion de molécules planes parmi celles que le modèle génère ?} Nous
y répondons par une démarche expérimentale en trois temps.

\begin{remarque}
Les valeurs chiffrées de cette section sont laissées en attente
(\tbd, \todo{}) : le corpus de résultats doit être régénéré sur la version
définitive de la table de voisinage avant intégration au mémoire.
\end{remarque}

\subsection{Protocole et banc d'essai}

Le banc d'essai utilise les squelettes benzénoïdes canoniques de la base
\emph{BenzaiDB}, de $h_3$ à $h_9$ hexagones (tableau~\ref{tab:corpus}). La
métrique principale est le \textbf{taux de planéité} : la proportion de
solutions jugées planes par le pipeline \xtb{}+ACP, complétée par le
\emph{nombre absolu} de molécules planes produites. Les campagnes sont
reproductibles, agrégées en base SQLite, et chaque configuration de contraintes
est comparée terme à terme aux autres sur le même corpus.

\begin{table}[h]
\centering
\begin{tabular}{lccccccc}
\toprule
$h$ & $h_3$ & $h_4$ & $h_5$ & $h_6$ & $h_7$ & $h_8$ & $h_9$ \\
\midrule
nombre de squelettes & \tbd & \tbd & \tbd & $44$ & $126$ & $376$ & $2\,418$ \\
\bottomrule
\end{tabular}
\caption{Nombre de squelettes benzénoïdes de BenzaiDB par taille $h$.}
\label{tab:corpus}
\end{table}

\subsection{v1 --- Modèle de base}

La première campagne mesure la référence : avec les seules contraintes
structurelles (section~\ref{sec:modele}), quelle proportion des solutions est
réellement plane après \xtb{} ? Le constat est que ce taux est médiocre et
\emph{décroît fortement} quand $h$ augmente (\tbd). Au-delà de la mesure, v1 a
surtout servi à mettre au point et stabiliser le pipeline complet
(reconstruction $\to$ MD $\to$ optimisation $\to$ ACP).

\subsection{v2 --- Contraintes globales : pentagones en bord}

L'idée de v2 est de traduire en contraintes \emph{a priori} deux régularités
empiriques observées sur les premières données : un équilibre entre pentagones
et heptagones, et la localisation des pentagones sur le bord. On introduit donc
deux contraintes optionnelles paramétrées :
\[
  \text{(C-SYM)}\ |n_5 - n_7| \le K_{\text{sym}},
  \qquad
  \text{(C-PB)}\ \sum_{v \in \partial \GD} \mathbf{1}[x_v = 5] \le K_{\text{pb}} .
\]
Les résultats sont nets. La contrainte de symétrie \textbf{C-SYM est sans effet}
mesurable : l'équilibre $n_5 = n_7$ est déjà imposé par C1, et le filtrage par la
table produit des compositions équilibrées. En revanche, \textbf{plafonner les
pentagones de bord est le levier dominant} : avec $K_{\text{pb}} = 1$ (préréglage
\texttt{pb1}), le taux de planéité gagne \todo{+X à +Y points} sur l'ensemble des
tailles, tandis que $K_{\text{pb}} = 0$ rend le problème insatisfiable. Le
préréglage \texttt{pb1} génère beaucoup moins de solutions, mais une proportion
bien plus élevée d'entre elles est plane (compromis volume / qualité favorable à
l'aval).

\begin{table}[h]
\centering
\begin{tabular}{lcc}
\toprule
$h$ & \texttt{baseline} (\% plan) & \texttt{pb1} (\% plan) \\
\midrule
$h_6$ & \tbd & \tbd \\
$h_7$ & \tbd & \tbd \\
$h_8$ & \tbd & \tbd \\
$h_9$ & \tbd & \tbd \\
\bottomrule
\end{tabular}
\caption{Taux de planéité : modèle de base contre préréglage \texttt{pb1}.}
\label{tab:v2}
\end{table}

\subsection{v3 --- Contraintes locales : appariement $5$--$7$ (Stone--Wales)}

v3 explore deux raffinements \emph{locaux}. Le premier borne la courbure de
Gauss--Bonnet dans un disque de rayon $R$ autour de chaque hexagone
($|n_5 - n_7| \le \tau$ localement, contrainte C-LC) ; il s'avère
\textbf{redondant} avec \texttt{pb1}. Le second impose l'\textbf{adjacence
$5$--$7$} (contrainte C5) : tout pentagone doit avoir un heptagone parmi ses
voisins immédiats, et réciproquement. C'est l'analogue discret du défaut
\emph{Stone--Wales} : un pentagone isolé porte une courbure $+\pi/3$ qui ne peut
être compensée qu'à distance ; en l'accolant à un heptagone ($-\pi/3$), la
courbure s'annule localement, condition nécessaire à un plongement plan.

La combinaison \texttt{pb1\_adj57} est la \textbf{configuration gagnante} : elle
améliore encore \texttt{pb1} et atteint un gain de \todo{+X à +Y points} par
rapport au modèle de base.

\begin{table}[h]
\centering
\begin{tabular}{lccc}
\toprule
$h$ & \texttt{baseline} & \texttt{pb1} & \texttt{pb1\_adj57} \\
\midrule
$h_6$ & \tbd & \tbd & \tbd \\
$h_7$ & \tbd & \tbd & \tbd \\
$h_8$ & \tbd & \tbd & \tbd \\
$h_9$ & \tbd & \tbd & \tbd \\
\bottomrule
\end{tabular}
\caption{Taux de planéité (\%) selon la configuration de contraintes.}
\label{tab:v3}
\end{table}

\subsection{Bilan expérimental}

L'étude désigne \texttt{pb1\_adj57} comme configuration recommandée (et par
défaut dans l'outil) : elle combine le plafond combinatoire des pentagones de
bord et l'appariement local $5$--$7$, et domine les autres réglages en taux de
planéité comme en coût de validation. Ses limites confirment l'analyse de la
section~\ref{sec:probleme} : sur les grands squelettes ($h_9$), le taux de
planéité s'effondre (\tbd), car les défauts de courbure s'accumulent au-delà de
ce que des contraintes locales peuvent contrôler --- ce qui rappelle la nature
intrinsèquement globale de la planéité.

% =====================================================================
\section{Conclusion et perspectives}
\label{sec:conclusion}
% =====================================================================

Ce stage a étendu aux cycles à $5$ et $7$ atomes le cadre, jusque-là limité aux
benzénoïdes, de génération de structures moléculaires par programmation par
contraintes. Nos contributions sont : un \textbf{modèle PPC} énumérant
exhaustivement les substitutions $5/6/7$ d'un benzénoïde ; une \textbf{table de
voisinage} sensible à la disposition des voisins et validée géométriquement, qui
sert de filtre chimique paramétrable ; un \textbf{pipeline de validation
physique} (reconstruction 3D, optimisation \xtb{} avec protocole de dynamique
moléculaire, test de planéité par ACP) ; une \textbf{étude expérimentale
systématique} qui identifie le couplage \texttt{pb1}\,+\,\texttt{adj-57} comme le
plus efficace pour améliorer le rendement de planéité, avec une interprétation en
termes de courbure discrète (motif Stone--Wales) ; et une \textbf{extension de
l'analyse électronique} (Kekulé, Clar, RBO) aux cycles $5$ et $7$.

Plusieurs perspectives se dégagent. La plus importante est de chercher à
\textbf{exprimer dans le modèle des contraintes globales de courbure}, afin de
réduire l'écart entre ce que les contraintes locales filtrent et la planéité
réelle. Sur le plan chimique, l'\textbf{analyse électronique} du corpus généré
mérite d'être approfondie (statistiques de Clar, RBO et radicaux, recherche de
corrélations), de même qu'un \textbf{RBO exact pour les cycles impairs} via le
théorème de Kasteleyn. Enfin, le \textbf{passage à l'échelle} au-delà de $h_9$ et
la poursuite du dialogue avec les chimistes --- pour affiner les critères
d'admissibilité et explorer la chimie radicalaire que les structures $5/7$
produisent naturellement --- prolongeraient naturellement ce travail.

% =====================================================================
%  BIBLIOGRAPHIE
%  (références à vérifier / compléter ; détails bibliographiques indicatifs)
% =====================================================================
\begin{thebibliography}{99}
\small

\bibitem{varet2022} A.~Varet,
\emph{Programmation par contraintes et chimie théorique : utilisation du
formalisme CSP pour résoudre des problématiques liées aux benzénoïdes},
thèse de doctorat, Aix-Marseille Université, 2022.

\bibitem{handbookcp} F.~Rossi, P.~van~Beek, T.~Walsh (éds.),
\emph{Handbook of Constraint Programming}, Elsevier, 2006.

\bibitem{lecoutre2009} C.~Lecoutre,
\emph{Constraint Networks: Techniques and Algorithms}, ISTE/Wiley, 2009.

\bibitem{pycsp3ace} C.~Lecoutre, N.~Szczepanski,
\emph{PyCSP3 : modélisation en programmation par contraintes} et le solveur
\emph{ACE} (Abscon Constraint Engine). \texttt{pycsp.org}.

\bibitem{benzai} Y.~Carissan, D.~Hagebaum-Reignier, N.~Prcovic, C.~Terrioux,
A.~Varet, \emph{BenzAI : un logiciel pour la conception et l'étude des
benzénoïdes}. \todo{référence exacte à compléter}

\bibitem{edmonds1965} J.~Edmonds,
\emph{Paths, trees, and flowers}, Canadian Journal of Mathematics, 17:449--467,
1965.

\bibitem{whitney} H.~Whitney,
\emph{Congruent graphs and the connectivity of graphs}, American Journal of
Mathematics, 54(1):150--168, 1932.

\bibitem{clar1972} E.~Clar,
\emph{The Aromatic Sextet}, John Wiley \& Sons, Londres, 1972.

\bibitem{huckel1931} E.~Hückel,
\emph{Quantentheoretische Beiträge zum Benzolproblem}, Zeitschrift für Physik,
70:204--286, 1931 (règle des $4n+2$ électrons $\pi$).

\bibitem{pauling} L.~Pauling,
\emph{The Nature of the Chemical Bond}, Cornell University Press, 3\textsuperscript{e}
édition, 1960.

\bibitem{randic} M.~Randić,
\emph{Aromaticity and conjugation}, Journal of the American Chemical Society,
99(2):444--450, 1977 (ordres de liaison et circuits conjugués).

\bibitem{rispoli2001} E.~Rispoli,
comptage des structures de Kekulé d'un benzénoïde par le déterminant de la
matrice de biadjacence, 2001. \todo{référence exacte à compléter}

\bibitem{kasteleyn1967} P.~W.~Kasteleyn,
\emph{Graph theory and crystal physics}, in \emph{Graph Theory and Theoretical
Physics}, Academic Press, p.~43--110, 1967.

\bibitem{stonewales} A.~J.~Stone, D.~J.~Wales,
\emph{Theoretical studies of icosahedral C$_{60}$ and some related species},
Chemical Physics Letters, 128(5--6):501--503, 1986.

\bibitem{xtb} C.~Bannwarth, S.~Ehlert, S.~Grimme,
\emph{GFN2-xTB — An Accurate and Broadly Parametrized Self-Consistent
Tight-Binding Quantum Chemical Method}, Journal of Chemical Theory and
Computation, 15(3):1652--1671, 2019.

\end{thebibliography}

\end{document}
